% §6 Further Results
%
% Condensed from old §6 (exactness) and §7 (extensions).
% Brief treatment — these support the JOSS paper, not this one.

%%====================================================================
\subsection{Algebraic exactness}
\label{sec:further-exactness}

The encoding pipeline tracks Pauli phases symbolically in the cyclic
group $\Zi = \{+1, -1, +i, -i\}$ rather than as floating-point complex
numbers.  Single-qubit Pauli multiplication is computed by exhaustive
pattern matching on the $4 \times 4$ Cayley table; no arithmetic is
involved.  The accumulated phase is folded into the coefficient by
negation and the swap $(a + bi) \mapsto (-b + ai)$, never by
\texttt{Complex.Multiply}.

\begin{theorem}[Algebraic exactness]
\label{thm:exactness}
Let $H = \sum_k c_k \prod_j O_j^{(k)}$ be a fermionic Hamiltonian
with coefficients $c_k \in \Qi$.  The Pauli-string representation
$\encode(H)$ satisfies: (i)~all intermediate phases are computed
exactly in $\Zi$; (ii)~the output coefficients lie in $\Qi$;
(iii)~term cancellation is exact (dictionary-keyed by Pauli
signature, not floating-point comparison).
\end{theorem}

\begin{proof}[Proof sketch]
Each pipeline stage — Majorana decomposition ($\pm\frac{1}{2}$,
$\pm\frac{i}{2}$), Pauli lookup ($\pm 1$, $\pm i$), qubit-by-qubit
multiplication ($\Zi$ accumulation), coefficient fold (negation and
swap) — introduces no rounding.  The only floating-point arithmetic
is addition of $\Qi$ coefficients in the like-term combination step,
which is exact for half-integers in IEEE~754.
\end{proof}

When input coefficients are not in $\Qi$ (e.g., from quantum chemistry
integrals), the pipeline still benefits: the only source of rounding is
the input data, not the encoding algebra.

%%====================================================================
\subsection{Extension to bosonic and mixed systems}
\label{sec:further-bosonic}

The normal-ordering algorithm is parameterized by a combining algebra
that specifies the fundamental swap relation:
\begin{align}
  \text{Fermionic (CAR):}\ & \an{j}\ad{j} = \delta_{jj} - \ad{j}\an{j},
  \label{eq:car-swap} \\
  \text{Bosonic (CCR):}\ & b_j b_j^\dagger = b_j^\dagger b_j + 1.
  \label{eq:ccr-swap}
\end{align}
The pipeline is identical for both statistics; only the swap rule
changes.  This abstraction is a concrete instance of the ``algebra as
parameter'' pattern.

Bosonic modes have unbounded occupation and must be truncated to a
maximum occupation $d$ before mapping to qubits.  We implement three
bosonic-to-qubit encodings: unary ($q = d$ qubits, weight~2), binary
($q = \lceil\log_2 d\rceil$ qubits, weight~$q$), and Gray-code
(same qubit count, improved average weight).  Unlike the fermionic
path, bosonic encodings introduce irrational coefficients
($\sqrt{n+1}$); algebraic exactness does not extend to this case.

Mixed fermion--boson Hamiltonians (electron--phonon, polariton--cavity)
are handled by sector tagging: each operator carries a
$\mathrm{Fermionic}$ or $\mathrm{Bosonic}$ label.  The normal-ordering
algorithm groups operators by sector, applies the appropriate combining
algebra within each block, and combines the results via Cartesian
product.  Cross-sector operators commute.  The qubit encoding tensors
the fermionic and bosonic registers on disjoint qubit sets.
